\section{The split space as a space of field strengths}
\label{sec:emfield}

The arena was introduced with an electromagnetic picture in mind: six directions, three of one sign and three of the other, paired
one to one like the components of a magnetic and an electric field. This section turns the picture into a theorem and measures its
price. The binding root $G$ of \cref{prop:twoK} is the key, and its role in the interior dynamics, where it is the operator that
makes bound motion possible, is established in \cref{sec:neutral}. Here it is used only as a linear map with $G^2=-\id$ and
$G=G^\dagger$.

\subsection{Arithmetic first}

The space $\Lambda^2\R^{1,3}$ with the induced metric $\langle\theta^a\wedge\theta^b,\theta^a\wedge\theta^b\rangle=\eta^{aa}\eta^{bb}$ is
negative on the three planes $\theta^0\wedge\theta^i$ and positive on the three planes $\theta^j\wedge\theta^k$. It therefore has
signature $(3,3)$, the Lorentz group acts on it by the adjoint representation, and the Hodge star satisfies $\star^2=-\id$. This is
counting, and it must be said before anything is built on it: every six-dimensional space of split signature is isometric to a
Lorentzian field-strength space, so exhibiting one isometry predicts nothing. What is not arithmetic is which subalgebra
$\so{1,3}\subset\so{3,3}$ realises the action, and whether the operator playing the Hodge star is one the arena already had.

\subsection{The Lorentz algebra of the binding root}

\begin{theorem}[The centraliser of $G$]
\label{thm:GLorentz}
The centraliser $\lG$ of $G$ in $\so{3,3}$ is isomorphic to $\so{3,\C}\cong\slr{2,\C}\cong\so{1,3}$. Its maximal compact
subalgebra is the $\so3$ of \emph{diagonal} rotations, each acting identically on a temporal plane $(e_i,e_j)$ and on its spatial
partner $(e_{i+3},e_{j+3})$.
\end{theorem}

\begin{proof}
Since $G^2=-\id$, the identification $z_j=x_j+ix_{j+3}$ makes $\R^6$ a complex space $\C^3$ on which $G$ acts as $i$, and
commuting with $G$ is $\C$-linearity. Under this identification
\[
\eta(X,X)=\sum_jx_j^2-\sum_jx_{j+3}^2=\Re q(z),\qquad q(z):=\sum_jz_j^2,
\]
with $q$ holomorphic. A $\C$-linear map $A$ preserving $\Re q$ preserves it on every phase rotation $e^{i\vartheta}z$, and
$q(e^{i\vartheta}z)=e^{2i\vartheta}q(z)$, so $\Re\bigl(e^{2i\vartheta}[q(Az)-q(z)]\bigr)=0$ for all $\vartheta$ forces $q(Az)=q(z)$. The
group is $O(3,\C)$ and its algebra $\so{3,\C}$, whose underlying real Lie algebra is $\slr{2,\C}\cong\so{1,3}$ with maximal compact
$\so3$. A rotation of $(e_i,e_j)$ alone does not commute with $G$, because $G$ carries $e_i$ to $e_{i+3}$; only the rotation
performed identically on the partner plane does.
\end{proof}

\begin{proposition}[Two Lorentz algebras, not conjugate]
\label{prop:notconj}
The interior algebra $\lint$ and $\lG$ are isomorphic as Lie algebras but not as subalgebras of $\so{3,3}$. The interior algebra
fixes the two-dimensional surplus plane pointwise, while $\lG$ contains simultaneous rotations of both triples and has no nonzero
common fixed vector. Conjugation preserves the dimension of the common fixed subspace, so the two are not conjugate in $GL(6,\R)$,
let alone in $O(3,3)$.
\end{proposition}

Write $L_i$ for the diagonal rotations and put, for cyclic $(i,j,k)$, $K_i=\tfrac12\bigl(e_ke_{j+3}-e_je_{k+3}\bigr)$. The six
operators $L_i,K_i$ span $\lG$.

\begin{theorem}[The field-strength reading]
\label{thm:emfield}
Write a vector of $\R^{3,3}$ as $X=\sum_i\bigl(B_ie_i+E_ie_{i+3}\bigr)$. Then
\[
[L_i,L_j]=\epsilon_{ijk}L_k,\qquad [L_i,K_j]=\epsilon_{ijk}K_k,\qquad [K_i,K_j]=-\epsilon_{ijk}L_k,
\]
both triples $\bB$ and $\bE$ transform as vectors under the $L_i$, and the $K_i$ act by
\begin{equation}
\label{eq:emboost}
\delta\bE=\bv\times\bB,\qquad \delta\bB=-\bv\times\bE .
\end{equation}
Identify $X$ with the two-form $F_X$ of components $(F_X)_{0i}=E_i$ and $(F_X)_{jk}=\epsilon_{ijk}B_i$. Then $G$ is the Hodge star,
$F_{GX}=\star F_X$ with $(\star F)_{ab}=\tfrac12\varepsilon_{abcd}F^{cd}$ and $\varepsilon_{0123}=+1$, and
\[
\eta(X,X)=|\bB|^2-|\bE|^2,\qquad \eta(X,GX)=-2\,\bE\cdot\bB
\]
are the only invariant quadratic forms. With $\lG$ as its Lorentz algebra, $\R^{3,3}$ \emph{is} $\Lambda^2\R^{1,3}$.
\end{theorem}

The orientation matters once, and we fix it here. With $F_{jk}=\epsilon_{ijk}B_i$ the vector $\bB$ is minus the textbook magnetic
field of the convention $F_{jk}=-\epsilon_{ijk}B_i$ \cite{Jackson}, so we write $\Btext=-\bB$, as announced in \cref{sec:conventions};
the electric field is the textbook one. Passing to a frame moving with velocity $\mathbf u$ changes the frame variables by
$\delta\bB=\mathbf u\times\bE$ and $\delta\bE=-\mathbf u\times\bB$, which is \eqref{eq:emboost} with $\bv=-\mathbf u$; in textbook
variables this is $\bE'=\bE+\mathbf u\times\Btext$ and $\Btext'=\Btext-\mathbf u\times\bE$, the familiar first-order transformation.

\begin{example}[A boost along the third axis]
\label{ex:emboost}
Take the pure magnetic field $X=e_1$, that is $\bB=(1,0,0)$ and $\bE=0$, and the boost generator $K_3=\tfrac12(e_2e_4-e_1e_5)$,
acting on vectors by the commutator. Since $e_1$ anticommutes with $e_2$ and with $e_4$ it commutes with $e_2e_4$, while
$[e_1e_5,e_1]=-2e_1^2e_5=-2e_5$; hence $[K_3,e_1]=e_5$. Under $\delta X=\delta\,[K_3,X]$ the field acquires $\bE=(0,\delta,0)$, in accordance
with $\delta\bE=\bv\times\bB$ for $\bv=(0,0,\delta)$. In the same way $[K_3,e_5]=e_1$, so a pure electric field $\bE=(0,1,0)$ acquires
$\bB=(\delta,0,0)=-\bv\times\bE$. The invariant $\eta(X,X)=|\bB|^2-|\bE|^2$ changes only at
second order, as it must.
\end{example}

\begin{corollary}[The invariant complex structure]
\label{cor:Gselected}
$\R^6$ is irreducible under $\lG$, and the commutant of $\lG$ in $\End(\R^6)$ is $\operatorname{span}\{\id,G\}\cong\C$. Hence, once
$\lG$ is taken as the Lorentz algebra, its invariant complex structure is $\pm G$ and nothing else. The interior algebra, by contrast,
has a commutant of dimension five, $\R$ on the four-plane and all of $M_2(\R)$ on the surplus plane.
\end{corollary}

The logic of this corollary runs from $G$ to $\lG$ and back to $\pm G$, so the choice of $G$ is recovered, not removed. It is a
consistency statement, a uniqueness conditional on taking $\lG$ as the Lorentz algebra. Removing the choice would require a
characterisation of $\lG$ that does not mention $G$, and we know none.

\subsection{The two readings exclude one another}

The price of the field-strength reading is not small.

\begin{proposition}[Exclusion]
\label{prop:emexclusive}
Under $\lint$ the six directions decompose as $V\oplus2\cdot\mathbf1$, a four-vector and two scalars; under $\lG$ they are
irreducible, with no fixed vector. The two representations are not similar, so a single theory cannot use both algebras as the
Lorentz action on the same six directions. If the vector space $\R^{3,3}$ is read as the electromagnetic field strength, those six
directions are not a four-plane plus a surplus pair, and the Dirac structure, the projector census, $P_*$ and the diagonal sector
all lose the arena on which they are built.
\end{proposition}

What is excluded is reading one six-dimensional space in two ways, not a theory that contains both kinds of module. A theory may
contain spacetime vectors and field strengths as different representations of one Lorentz algebra, which is the ordinary
arrangement, and $\Cl(3,3)$ realises it with the vectors at grade one and the field strength at grade two.

\begin{proposition}[The field strength at grade two]
\label{prop:bivectorF}
Let $\Lambda^2E$ be the span of the six bivectors of the interior four-plane. It is closed under the adjoint action of $\lint$, on
which it is the adjoint representation; the trace form $\tfrac18\Tr(XY)$ on it has signature $(3,3)$; and the interior volume
$\Omega=\gamma^0\gamma^1\gamma^2\gamma^3$, with $\Omega^2=-\id$, commutes with $\lint$, preserves $\Lambda^2E$ by left multiplication,
exchanges $\gamma^0\gamma^i$ with the spatial bivectors and acts on $\Lambda^2E$ as the Hodge star.
\end{proposition}

Grade three behaves differently under the two algebras, and the difference sharpens the exclusion. Under $\lG$ one has
$(\Lambda^3\R^6)_\C=2\,(0,0)\oplus2\,(1,1)$, of dimension $2+18$, so no field-strength module $(1,0)\oplus(0,1)$ occurs at grade three
at all: trivectors cannot carry the field for this action. Under $\lint$, $\R^6=E\oplus\R^2$ gives $\Lambda^3\R^6\cong2E\oplus2\Lambda^2E$,
which contains two field-strength modules. The quadratic Casimir distinguishes the cases cleanly, with spectrum $\{0^{(2)},8^{(18)}\}$
on $\Lambda^3\R^6$ under $\lG$ and $\{3^{(8)},4^{(12)}\}$ under $\lint$, the value $4$ being that of $(1,0)\oplus(0,1)$ and $3$ that of a
four-vector. An explicit attempt confirms it: the six trivectors $C_i=\sum_{j,k}\epsilon_{ijk}e_jN_k$, $N_j=e_{k+3}e_{l+3}$ cyclically,
together with $D_i=\mathcal VC_i$, transform as two vectors under the $L_i$, but all eighteen brackets $[K_i,C_j]$, $[K_i,D_j]$ leave
their span, for instance $\tfrac12[K_1,C_1]=e_2e_3e_4+e_4e_5e_6$, and the Lorentz closure of the span has dimension $18$.

\subsection{How much of Maxwell's theory is forced}
\label{sec:maxwell}

\Cref{thm:emfield} supplies a field with the right transformation law and nothing that the field satisfies. How much of the
dynamics does the transformation law fix by itself? The answer is: the form of the equations completely, and nothing else.

\begin{theorem}[Exactly two first-order covariant operators]
\label{thm:twoops}
Consider linear first-order differential operators with constant coefficients on Minkowski space, acting on a field $F$ valued in
$\Lambda^2\R^{1,3}$ and taking values in the vector representation. The derivative $\partial F$ lies in $\Lambda^2\otimes V$, of
dimension $24$, and under $\so{1,3}$
\[
\Lambda^2\otimes V=\underbrace{(\tfrac32,\tfrac12)\oplus(\tfrac12,\tfrac32)}_{16}\ \oplus\ \underbrace{2\,(\tfrac12,\tfrac12)}_{8},
\]
so the vector module occurs with multiplicity exactly two, and no other module of dimension four occurs. The two copies are realised
by
\[
(dF)_{\mu\nu\rho}=3\,\partial_{[\mu}F_{\nu\rho]}\quad(\text{through }\Lambda^3\cong V),\qquad (\delta F)_\nu=\partial^\mu F_{\mu\nu},
\]
and every covariant operator of the class is $\mathcal DF=a\,\delta F+b\,\delta(\star F)$, since $\delta\star=\pm\star d$.
\end{theorem}

The Casimir operator confirms the count: on $\Lambda^2\otimes V$ its spectrum is $\{3^{(8)},9^{(16)}\}$, and the value $3$ is the
vector one. Nothing in the theorem constrains operators of higher order, operators with variable coefficients, or operators valued
in the remaining $(\tfrac32,\tfrac12)\oplus(\tfrac12,\tfrac32)$.

\begin{corollary}[The form of Maxwell's equations]
\label{cor:maxwellform}
Within this class the two available equations are $dF=\mathcal J_m$ and $d\star F=\mathcal J_e$. Covariance does not rule out a
magnetic current, which is a covariant field exactly as an electric current is; it identifies the two operators and leaves open
which equations are imposed and which sources are admitted. The source-free system $dF=\delta F=0$ is the only choice that needs no
external data, and within this linear class it is the only one invariant under duality, because the Hodge star $G$ exchanges the two
operators. Finally, on a contractible region $dF=0$ gives $F=dA$ by the Poincar\'e lemma, so the gauge potential is a consequence of
the first equation rather than an assumption.
\end{corollary}

In the variables of the frame the system takes a compact form, which we use repeatedly below.

\begin{proposition}[Maxwell's equations on the frame]
\label{prop:RS}
Read $\R^{3,3}$ as $\C^3$ through $G$, $z_j=B_j+iE_j$. Then $dF=0$ and $d\star F=0$ for the two-form $F_X$ of
\cref{thm:emfield} are
\begin{equation}
\label{eq:RS}
\partial_tX=-G\,(\curl X),\qquad \nabla\cdot X=0,
\end{equation}
that is $\partial_tz=-i\,\curl z$ and $\nabla\cdot z=0$, with the curl built from the rotation generators $L_i$. Splitting into the
two triples gives $\partial_t\bB=\curl\bE$, $\partial_t\bE=-\curl\bB$ and $\nabla\cdot\bE=\nabla\cdot\bB=0$; in textbook variables these
are Faraday's law $\partial_t\Btext=-\curl\bE$ and the Amp\`ere--Maxwell law $\partial_t\bE=\curl\Btext$ in vacuum. With a source,
$\partial_\mu F^{\mu\nu}=(\rho,\bj)$ reads
\begin{equation}
\label{eq:RSsource}
\partial_tz=-i\,\curl z-i\,\bj,\qquad \nabla\cdot z=i\rho ,
\end{equation}
and taking the divergence of the first equation gives $\partial_t\rho+\nabla\cdot\bj=0$: continuity is a consequence of the system and
the compatibility condition that any proposed current must satisfy.
\end{proposition}

Up to a factor $i$, $z$ is the Riemann--Silberstein vector $\bE+i\Btext$, which obeys
$i\partial_t(\bE+i\Btext)=\curl(\bE+i\Btext)$ \cite{Silberstein1907,BialynickiBirula13}, and the holomorphic form of the proof of
\cref{thm:GLorentz} has
\[
\Re q(z)=|\bB|^2-|\bE|^2,\qquad \Im q(z)=2\,\bE\cdot\bB ,
\]
so the two Maxwell invariants are the real and imaginary parts of one holomorphic quadratic form. Three classical facts fall into
place. The duality rotation is the phase rotation $z\mapsto e^{i\vartheta}z$ that the proof of \cref{thm:GLorentz} already used, and
it turns the pair of invariants by $2\vartheta$. The two quadratic forms that will compete in the interior dynamics of
\cref{sec:shifted,sec:neutral} reappear with physical names: the $\eta$-form $\Re q$ is the Lorentz invariant, while the Hermitian
form $|z|^2=|\bB|^2+|\bE|^2$ is twice the energy density and depends on the frame. And the remaining sesquilinear covariant is the
Poynting vector, $\tfrac i2\,\bar z\times z=\bE\times\bB$, so that $\bE\times\Btext=-\tfrac i2\,\bar z\times z$.

\begin{example}[A circularly polarised wave on the frame]
\label{ex:circular}
The field $z=(1,i,0)\,e^{ik(x_3-t)}$ solves \eqref{eq:RS}: $\nabla\cdot z=0$ because $z$ has no third component, and
$\curl z=ik\,(-i,1,0)e^{ik(x_3-t)}=k\,(1,i,0)e^{ik(x_3-t)}$, so $-i\curl z=-ik\,z=\partial_tz$. Its invariant is
$q(z)=(1+i^2)e^{2ik(x_3-t)}=0$, so $|\bB|=|\bE|$ and $\bE\cdot\bB=0$ at every point: the wave is null. Its Poynting vector
$\tfrac i2\bar z\times z=\tfrac i2(1,-i,0)\times(1,i,0)=\tfrac i2(0,0,2i)=(0,0,-1)$ is $\bE\times\bB$, so $\bE\times\Btext=(0,0,1)$
points along the direction of propagation, as it should.
\end{example}

\subsection{Frame vectors as spacetime planes}

Each pairing plane $(e_i,e_{i+3})$ carries one complex component $z_i=B_i+iE_i$, and a frame vector is at the same time a bivector of
spacetime. The two descriptions meet in the following way.

\begin{proposition}[Frame vectors as planes]
\label{prop:planes}
Let $X=\sum_i(B_ie_i+E_ie_{i+3})$ and $F_X$ as above. Then $\Pf F_X=\bE\cdot\bB$ and $\det F_X=(\bE\cdot\bB)^2$. A nonzero $F_X$ is a
single plane $u\wedge v$ exactly when $\bE\cdot\bB=0$, that is when $\Im q=0$. Otherwise there is a Lorentz frame in which
$F_X=a\,\theta^0\wedge\theta^1+b\,\theta^2\wedge\theta^3$ is a pair of orthogonal planes, and the duality rotation $e^{-\vartheta G}$ with
$\tan\vartheta=a/b$ carries it to the single plane $\sqrt{a^2+b^2}\,\theta^2\wedge\theta^3$, the complexion of Rainich and of Misner
and Wheeler \cite{Rainich1925,MisnerWheeler57}.
\end{proposition}

\Cref{fig:qplane} shows the content of this proposition in the plane of the complex invariant. The real axis is the locus of single
planes, the positive half of magnetic type and the negative half of electric type, the origin is the null cone, and the duality
rotation moves a field around a circle at twice its own angle. In \cref{sec:neutral} the bound orbits of the interior dynamics will
turn out to be exactly such rotations.

\begin{figure}[t]
\centering
\begin{tikzpicture}[scale=1.15,>=Stealth]
\draw[->,gray] (-3.0,0) -- (3.2,0) node[right] {\small $\Re q=|\bB|^2-|\bE|^2$};
\draw[->,gray] (0,-2.6) -- (0,2.8) node[above] {\small $\Im q=2\,\bE\cdot\bB$};
\draw[scol,line width=2.2pt] (-2.9,0) -- (-0.08,0);
\draw[tcol,line width=2.2pt] (0.08,0) -- (2.9,0);
\node[tcol,anchor=north west] at (0.15,-0.08) {\small single plane, magnetic};
\node[scol,anchor=north east] at (-0.15,-0.08) {\small single plane, electric};
\fill[black] (0,0) circle (2pt);
\node[anchor=north west] at (0.05,0.45) {\small null};
\draw[gcol,thick,dashed] (0,0) circle (1.8);
\coordinate (q0) at ({1.8*cos(35)},{1.8*sin(35)});
\fill[gcol] (q0) circle (2.3pt);
\node[gcol,right=3pt] at (q0) {\small $q_0$};
\draw[->,gcol,thick] ({2.2*cos(35)},{2.2*sin(35)}) arc (35:95:2.2);
\node[gcol,align=left,anchor=west] at (1.75,2.05) {\small duality $e^{\vartheta G}$:\\ \small $q\mapsto e^{2i\vartheta}q$};
\foreach \a in {0,180}{\fill[ocol] ({1.8*cos(\a)},{1.8*sin(\a)}) circle (2.0pt);}
\node[ocol,anchor=north east,align=right] at (-1.0,-1.55) {\small each crossing of the real axis\\ \small is met twice per turn of $\vartheta$};
\end{tikzpicture}
\caption{The plane of the complex invariant $q=z\cdot z$ of a frame vector (\cref{thm:emfield,prop:planes}). The real axis is the
locus $\bE\cdot\bB=0$ on which the spacetime two-form is a single plane, magnetic on the right and electric on the left; the origin is
the null cone. A duality rotation $e^{\vartheta G}$ preserves $|q|$ and turns $q$ at twice the rate, so a non-null field becomes a single
plane four times as $\vartheta$ runs once around.}
\label{fig:qplane}
\end{figure}

Everything in this section is a statement about a representation and its invariants. \Cref{thm:twoops} is a multiplicity count that
fixes the form of the equations and nothing about their content: it does not produce a source, does not say that a source is the
current of anything, supplies no coupling constant and no Lagrangian. Nor is any of it special to this arena, since the count uses only
$\so{1,3}$ and $\Lambda^2$. What the arena adds is that the objects the count needs, the module, the duality and the invariants, were
already present, so that the system can be written in the arena's own variables \eqref{eq:RS}. The deflationary test of
\cref{sec:deflation} applies to it in full.
